\section*{Chapter \ref{ch:g5:series-parallel-circuits} --- Series and Parallel Circuits}

\begin{solution}{exo:g5:series-parallel-circuits:1}
A \omterm{def:g5:series-parallel-circuits:parallel}{branch} is one of the separate roads the circuit forks into; a
\omterm{def:g5:series-parallel-circuits:parallel}{junction} is a point where roads fork or rejoin. Two \omterm{def:g3:first-electric-circuit:bulb}{bulbs} \omterm{def:g5:series-parallel-circuits:parallel}{in
parallel} ride two bulb-carrying \omterm{def:g5:series-parallel-circuits:parallel}{branches}.
\end{solution}

\begin{solution}{exo:g5:series-parallel-circuits:2}
Independence (a gap stops only its own \omterm{def:g5:series-parallel-circuits:parallel}{branch}); full brightness
(each \omterm{def:g3:first-electric-circuit:bulb}{bulb} shines as if alone, the \omterm{def:g3:first-electric-circuit:battery}{battery} draining faster in
return); command posts (a \omterm{def:g5:series-parallel-circuits:parallel}{branch} \omterm{ex:g3:first-electric-circuit:switch}{switch} rules its \omterm{def:g5:series-parallel-circuits:parallel}{branch}, a \omterm{ex:g3:first-electric-circuit:switch}{switch}
before the fork rules all).
\end{solution}

\begin{solution}{exo:g5:series-parallel-circuits:3}
\omterm{def:g3:first-electric-circuit:bulb}{Bulb} 2 burns on, undisturbed --- its own road is complete. \omterm{def:g4:series-circuits:series}{In
series} it went dark, because there the two \omterm{def:g3:first-electric-circuit:bulb}{bulbs} shared one single
road, and any gap broke it for both.
\end{solution}

\begin{solution}{exo:g5:series-parallel-circuits:4}
Pair B (parallel) shines brighter per \omterm{def:g3:first-electric-circuit:bulb}{bulb} --- each at full
brightness. Pair B also drains its \omterm{def:g3:first-electric-circuit:battery}{battery} faster: it feeds two
full \omterm{def:g5:series-parallel-circuits:parallel}{branches} at once.
\end{solution}

\begin{solution}{exo:g5:series-parallel-circuits:5}
Both closed: both \omterm{def:g3:first-electric-circuit:bulb}{bulbs} shine. Main open: both dark --- it rules
the common road. Main closed, \omterm{def:g5:series-parallel-circuits:parallel}{branch} \omterm{ex:g3:first-electric-circuit:switch}{switch} open: \omterm{def:g3:first-electric-circuit:bulb}{bulb} 1 shines,
\omterm{def:g3:first-electric-circuit:bulb}{bulb} 2 dark.
\end{solution}

\begin{solution}{exo:g5:series-parallel-circuits:6}
\omterm{def:g4:series-circuits:series}{In series}, one dead \omterm{def:g3:first-electric-circuit:bulb}{bulb} (or one switched-off lamp!) would darken
the whole house, and everything would dim as more \omterm{def:g1:light-and-shadows:source}{lights} joined.
Comforts from parallel wiring: each room's \omterm{ex:g3:first-electric-circuit:switch}{switch} rules only its
room; a burnt-out kitchen \omterm{def:g3:first-electric-circuit:bulb}{bulb} leaves the rest of the house lit;
every lamp burns at full brightness.
\end{solution}

\begin{solution}{exo:g5:series-parallel-circuits:7}
Before the fork --- on the common road where the current enters the
home: the main \omterm{ex:g3:first-electric-circuit:switch}{switch} in the fuse box. It cuts the whole house at
once, used for repairs and emergencies.
\end{solution}

\begin{solution}{exo:g5:series-parallel-circuits:8}
\omterm{def:g5:series-parallel-circuits:parallel}{In parallel}: switching the ceiling \omterm{def:g1:light-and-shadows:source}{light} off does not kill the
bedside lamp, and unplugging the lamp does not darken the ceiling
--- each rides its own \omterm{def:g5:series-parallel-circuits:parallel}{branch}.
\end{solution}

\begin{solution}{exo:g5:series-parallel-circuits:9}
With our plain \omterm{ex:g3:first-electric-circuit:switch}{switches} the designs keep failing: a \omterm{ex:g3:first-electric-circuit:switch}{switch} \omterm{def:g4:series-circuits:series}{in
series} with the \omterm{def:g3:first-electric-circuit:bulb}{bulb} can be opened from one end only; two plain
\omterm{ex:g3:first-electric-circuit:switch}{switches} \omterm{def:g4:series-circuits:series}{in series} demand \emph{both} closed (either end can turn
the \omterm{def:g1:light-and-shadows:source}{light} off but not always on); two \omterm{def:g5:series-parallel-circuits:parallel}{in parallel} demand both open
for darkness. The hallway comfort needs a \omterm{ex:g3:first-electric-circuit:switch}{switch} that \emph{chooses
between two roads} instead of just breaking one --- the two-way
\omterm{ex:g3:first-electric-circuit:switch}{switch}, which real hallways use in matched pairs.
\end{solution}

\begin{solution}{exo:g5:series-parallel-circuits:10}
The comfort: the \omterm{def:g3:first-electric-circuit:bulb}{bulbs} ride parallel \omterm{def:g5:series-parallel-circuits:parallel}{branches}, so a dead \omterm{def:g3:first-electric-circuit:bulb}{bulb}
darkens only itself. The warning: every missing \omterm{def:g3:first-electric-circuit:bulb}{bulb} removes a
\omterm{def:g5:series-parallel-circuits:parallel}{branch}, and the \omterm{def:g3:first-electric-circuit:battery}{battery} (or the string's little power supply)
pushes its full effort through the \omterm{def:g5:series-parallel-circuits:parallel}{branches} that remain --- with
too many gone, the survivors each receive more than they were
built for, burn hotter, and fail young.
\end{solution}

\begin{solution}{exo:g5:series-parallel-circuits:11}
Onto the \omterm{ex:g3:first-electric-circuit:switch}{branch-switch}'s own stretch of road: between the \omterm{def:g5:series-parallel-circuits:parallel}{branch}
\omterm{ex:g3:first-electric-circuit:switch}{switch} and \omterm{def:g3:first-electric-circuit:bulb}{bulb} 2 (before their \omterm{def:g5:series-parallel-circuits:parallel}{branch} rejoins the rest). There it
\omterm{def:g1:light-and-shadows:source}{lights} exactly when that \omterm{def:g5:series-parallel-circuits:parallel}{branch} runs --- and it sits \omterm{def:g4:series-circuits:series}{in series} with
\omterm{def:g3:first-electric-circuit:bulb}{bulb} 2, both riding the same \omterm{def:g5:series-parallel-circuits:parallel}{branch}.
\end{solution}
